AL ver las figuras previas, podemos demostrar la utilidad de Redis para delegación de tareas, haciendo más eficiente cada uno de los procesos a realizar.

\Chapter{discusión y conclusión}

Finalmente, vemos que mientras una caché como la B se actualiza, en A si bien, ya también está almacenado pero con otra key dado a la función de preload. Entonces, si queremos ver inmediatamente reflejados los datos que tenemos almacenados, basta con solo acceder a students, ya que en cada update o delete precargará información con está clave. Esto permite que no se tenga que estar accediendo siempre a mongo y la consulta se realiza más rápidamente.

Actualmente los nodos solo fungen como delegadores de trabajo, gracias a nuestros middleware, sin embargo, todavía podrían hacer procesos más avanzados en bases de datos distribuidas. L ejemplo que previamnet ese dió fue el de los workers para la inserción asíncorna de los datos.

Además, se podrían utilizar más tecnologías como redis sentinel. Con él se puede evitar el fallback manual simple que realizamos, que una forma de como lo \textit{preveníamos} era con un acceso aleatorio, sin embargo para evitar hacer esto, redis sentinel hace un análisis de cual nodo debería de utilizar y le ventar a esclado en caso de fallback. Además de que se encarga del balanceo de cargas. Redis clúster nos ayudaría a evitar construir manualmente en el docker compose cada una de las replicas que se crean. Se debe e recordar que se levantan 8 servicios.

El diseño de nuestro backend o el diseño de nuestra red de conexión dependerá mucho de lo que buscamos realizar, que problemas queremos evitar y propiciar a la escabilidad del software. Los patrones de diseño son importantes para este tipo de proyectos y son fundamentales para el diseño y desarrollo de software con buenas prácticas. La teoría revisada en el marco téorico es fundamental y necesaria para saber que algoritmo o patrón se puede implementar y cual podría ayudar para nuestra problemática
